\documentclass[11pt]{article}
\usepackage[margin=1in]{geometry}
\usepackage{fontspec}
\usepackage{amsmath,amssymb}
\usepackage{mathrsfs}
\usepackage{amsthm}
\newtheorem{proposition}{Proposition}
\usepackage{enumitem}
\usepackage{hyperref}
\usepackage{titlesec}
\usepackage{parskip}
\DeclareMathOperator*{\argmin}{arg\,min}
\hypersetup{colorlinks=true,linkcolor=black,citecolor=black,urlcolor=black}

\titleformat{\section}{\normalfont\Large\bfseries}{\thesection}{1em}{}
\titleformat{\subsection}{\normalfont\large\bfseries}{\thesubsection}{1em}{}

\title{\textbf{Sufficiency Under Compression: Exact Decoding, Admissible Projection, and Operational Coupling in Clio and Hydra}}
\author{Flyxion \\ \small Independent Researcher}
\date{August 2026}

\begin{document}
\maketitle

\begin{abstract}
Compression is often described by how much smaller a representation becomes, as though size reduction were the phenomenon of interest. This essay argues that the governing question is different: which distinctions may be removed, combined, or made only indirectly computable without preventing a downstream system from performing its intended operation. Three architectures are examined as points on a single spectrum answering that question differently. A JPEG Huffman-decoding patent replaces an exhaustive $2^{16}$-entry lookup table with a compact intensional construction, and does so with exact, finitely verifiable correctness for the specific decoding task it targets. Clio's projection collapses a state to an equivalence class under an admissibility relation. This essay reconciles three formalizations of that construction found across the wider research program, a fixed admissibility quotient, a convergent iterative estimator of that quotient, and a superseded co-evolving variant, and shows exactly when and how badly such a projection can fail using a quantitative collapse observable. It then complicates the convergence result itself with a recorded case where non-convergence was diagnostic of an inadmissible quotient rather than defective estimation. Hydra's transition operator combines perception, relevance, and memory signals in a way that may be mere joint retrieval or may be genuine operational coupling. The essay develops a non-vacuous separability test to tell the two apart, applies it to Hydra's actual named layer equations to locate where coupling might be hidden in shared parameters rather than top-level combination rules, and adds a categorical gluing-obstruction diagnostic alongside a residue-field complication affecting the test's own intervention protocol and a mixture-of-experts contrast case that looks coupled but tests separable. The unifying formal notion throughout is $F$-sufficiency: a quotient retains everything a declared function or function class requires, and nothing more is guaranteed. The paper closes by composing the two case studies into a single pipeline, from quotient to gluing to transition, and showing that repair has two structurally distinct entry points depending on which stage's guarantee actually failed.
\end{abstract}

\section{Introduction: Compression versus Size Reduction}

A representation can be made smaller in many ways, and most of them are uninteresting on their own. The interesting question is not how much smaller a representation has become but which distinctions in the original object it still preserves, in what sense it preserves them, and for whom. A representation that discards a distinction no downstream operation will ever need has lost nothing. A representation that discards a distinction some future operation silently depends on has incurred a debt that will eventually come due, whether or not the representation's designer intended to borrow against it.

This essay examines three architectures that answer this question in structurally different ways, ordered by how strong a guarantee each provides. A patented JPEG Huffman-decoding method replaces a fully enumerated lookup table with a compact computed procedure, and does so with a correctness guarantee that can be checked exhaustively because the input alphabet is finite. Clio, a component of the author's broader admissibility-theoretic research program, projects a state onto its normalized direction, discarding magnitude on the working assumption that magnitude is irrelevant to whatever operation consumes the projection; this assumption is not verified by the projection itself and can fail. Hydra, a cognitive architecture combining perception, relevance activation, and memory signals within a shared field-theoretic substrate, raises a question the first two cases do not: whether its component signals are merely retrieved together or whether they participate in the system's transition dynamics in a way that cannot be decomposed into independent contributions at all.

These are not three unrelated case studies assembled for variety. They are three instances of a single formal pattern, a map from an object to a smaller representation and then to an output, together with three distinct strengths of guarantee that the smaller representation was sufficient for the output to be correct. Section 2 states this pattern formally. Sections 3 through 5 work through the three cases in order of decreasing guarantee strength. Section 6 states the resulting hierarchy precisely. Section 7 connects the whole apparatus to the admissibility and repair vocabulary used elsewhere in this research program: compression creates an obligation, and repair is what happens when that obligation comes due.

\section{Formal Setup: $F$-Sufficiency}

Let $x$ be an object drawn from a domain $X$, let $q: X \to Q$ be a quotient map into a smaller representation space, and let $g: Q \to Y$ be a downstream operation applied to the quotient. Write $\widehat{f} = g \circ q$ for the composite, intended to approximate or reproduce some target function $f: X \to Y$:
\[
x \xrightarrow{\;q\;} q(x) \xrightarrow{\;g\;} \widehat{f}(x).
\]
The question this essay asks of every architecture it examines is the same: under what conditions does $\widehat{f}$ agree with $f$, and how strong is the argument for that agreement.

This is a deterministic, functional relative of Fisher--Neyman sufficiency in classical statistics, where a statistic $T(X)$ is sufficient for a parameter $\theta$ if the conditional distribution of $X$ given $T(X)$ does not depend on $\theta$, meaning $T$ retains everything the sample carries about $\theta$ and nothing is lost by discarding the rest. The present setting has no parameter and no likelihood family; what varies here is not an unknown parameter but a declared downstream function or function class the quotient is required to serve. To avoid quietly importing statistical assumptions that do not hold here, this essay calls the relevant property \emph{$F$-sufficiency} rather than Fisher--Neyman sufficiency, and defines it directly in terms of the function or function class of interest.

For a single function $f$, the quotient $q$ is $F$-sufficient for $f$ if
\[
q(x) = q(y) \implies f(x) = f(y)
\]
for all $x, y \in X$: whenever the quotient cannot distinguish two inputs, neither can the target function, so nothing $f$ needed has been discarded. For a function class $\mathcal{F}$, the quotient is $F$-sufficient for $\mathcal{F}$ if
\[
q(x) = q(y) \implies \forall f \in \mathcal{F},\; f(x) = f(y).
\]
This class version is the one that matters for Clio in Section 4: a projection is not simply sufficient or insufficient in the abstract, it is sufficient \emph{relative to} whichever function class the system's designer expects future operations to belong to, and a later operation lying outside that class is precisely a case where the projection's sufficiency claim fails.

Three distinguishable strengths of $F$-sufficiency claim recur across the cases below. \emph{Exact sufficiency} requires $g(q(x)) = f(x)$ for every $x$ in the domain, with no approximation. \emph{Task-relative sufficiency} permits an explicitly bounded loss under a specified task and input distribution, of the form $\mathbb{E}[L(f(X), g(q(X)))] \leq \varepsilon$ for a stated loss $L$. \emph{Operational coupling}, treated separately in Section 5, is not a sufficiency claim about reconstruction accuracy at all; it is a claim about whether the components of a joint object participate \emph{separably} in a downstream transition, and requires its own formal treatment because the vacuous version of a separability claim is trivially satisfiable, as Section 5 shows directly.

\section{Exact Functional Sufficiency: Huffman Length Classification}

The clean baseline for this hierarchy is a method described in U.S. Patent 6,373,412 B1 (Mitchell, Cazes, and Leeder, assigned to IBM, filed 2000, granted 2002), covering fast JPEG Huffman encoding and decoding \cite{mitchell2002}. The baseline JPEG decoding problem is: given up to sixteen bits of compressed data, determine the length $n$ of the valid Huffman code those bits begin with, and return the symbol that code represents. The naive solution is extensional enumeration: a lookup table with $2^{16}$ entries, one for every possible 16-bit prefix, each entry storing the corresponding symbol directly. This table is exact by construction, since it is built by writing down the answer for every input, but it is enormous relative to the actual information content of a Huffman table, which for a baseline JPEG image contains at most a few hundred defined codes.

The patent's construction replaces this table with an intensional representation. For each candidate length $n$ from 1 to 16, it precomputes a threshold value $\mathrm{max}(n)$, the numeric value one greater than the largest $n$-bit code actually used at that length, and an $\mathrm{offset}(n)$, used to translate a valid $n$-bit code into an index into a much smaller table of only the symbols actually present. Decoding proceeds by testing the incoming bit pattern against $\mathrm{max}(1), \mathrm{max}(2), \ldots$ in increasing order until the first length at which the pattern's leading bits are less than the threshold; this length is then combined with $\mathrm{offset}(n)$ to compute a direct index into the compact symbol table.

Let $q_{\mathrm{JPEG}}$ denote the map from a 16-bit input pattern to the pair (code length $n$, computed index), and let $g_{\mathrm{JPEG}}$ denote the lookup into the compact symbol table using that pair. The patent's central technical claim, in the vocabulary of Section 2, is that
\[
g_{\mathrm{JPEG}}(q_{\mathrm{JPEG}}(x)) = f_{\mathrm{JPEG}}(x)
\]
for every valid input $x$, where $f_{\mathrm{JPEG}}$ is the extensional $2^{16}$-entry table's own output. This is exact sufficiency, and it is a stronger claim than exact sufficiency typically permits elsewhere in this essay, because the domain here is finite and completely specified: the equivalence between the intensional and extensional constructions can be verified by exhaustive comparison against every one of the $2^{16}$ possible inputs. This does not make JPEG the only case of exact sufficiency in general; exact functional sufficiency can hold over infinite domains as well, wherever a quotient and a target function happen to align exactly. What is distinctive about the JPEG case is narrower and worth stating precisely: it is the paper's only case of \emph{finitely verifiable} exact sufficiency, where the claim can be checked by brute enumeration rather than argued from other premises. Sections 4 and 5 will not have this luxury.

\subsection{The Epistemic Status of Exact Sufficiency}

It is worth isolating exactly what the finite verification above does and does not establish, since it is easy to overstate. Every possible input in the 16-bit prefix space can be checked against the specified decoder's output, so no input falls outside the guarantee for that one fixed task, $f_{\mathrm{JPEG}}$. This is not a guarantee against every conceivable future function over the same bits. A different downstream task, recovering the original bit prefix before truncation, characterizing a corruption pattern, or auditing the statistical distribution of code lengths actually used, would constitute a different sufficiency claim, one this section's exhaustive verification of $f_{\mathrm{JPEG}}$ says nothing about, and the $\mathrm{max}(n)/\mathrm{offset}(n)$ quotient may well be insufficient for it. Exactness here is total relative to one declared function and finite relative to one declared domain; it is not total relative to every function that domain could support.

\section{Admissibility-Relative Sufficiency: Clio}

Clio, as used elsewhere in this research program, projects a vector $x$ onto its $L_2$-normalized direction:
\[
q_{\mathrm{Clio}}(x) = \frac{x}{\lVert x \rVert_2}.
\]
This projection identifies all positive scalar multiples of a given vector:
\[
x \sim_{\mathrm{Clio}} y \iff y = \alpha x, \quad \alpha > 0,
\]
discarding magnitude entirely and retaining only direction. Unlike the JPEG case, this projection's correctness is not something the construction itself establishes. Whether $q_{\mathrm{Clio}}$ is $F$-sufficient for a given downstream function $f$ depends entirely on whether $f$ happens to be invariant under the equivalence relation it induces. The decisive condition, stated exactly as in Section 2, is
\[
x \sim_{\mathrm{Clio}} y \implies f(x) = f(y),
\]
or, where exact invariance is too strong a demand, its task-relative relaxation
\[
\mathbb{E}\left[ L\bigl(f(X), g(q_{\mathrm{Clio}}(X))\bigr) \right] \leq \varepsilon
\]
for a specified loss $L$ and input distribution.

This condition is a claim about $f$, not a property $q_{\mathrm{Clio}}$ can certify on its own. It is entirely possible to construct a task for which the projection fails outright. Consider a function $f_{\mathrm{mag}}(x) = \mathbb{1}[\lVert x \rVert_2 > \tau]$ for some threshold $\tau$, a function whose entire content is a judgment about magnitude. For any two vectors $x$ and $y = \alpha x$ with $\alpha$ chosen so that exactly one of $x, y$ crosses the threshold $\tau$, $x \sim_{\mathrm{Clio}} y$ holds by construction while $f_{\mathrm{mag}}(x) \neq f_{\mathrm{mag}}(y)$. The projection is not merely imperfect here; it is maximally uninformative, since $q_{\mathrm{Clio}}$ carries no information whatsoever about which side of $\tau$ a given vector falls on. This is not a pathological edge case invented to embarrass the construction. It is the generic situation whenever a task's relevant signal is carried by scale rather than by direction, and it is exactly the kind of task a system built around Clio would fail silently on unless the invariance assumption were checked rather than assumed.

The role of this example is to prevent a specific error: treating the act of projecting as though it were itself evidence that the discarded dimension was irrelevant. It is not. Sufficiency under $q_{\mathrm{Clio}}$ is a claim about which function class $\mathcal{F}$ future downstream operations are expected to belong to, in the sense of Section 2's class version of $F$-sufficiency, and that claim can be true of the operations a system was designed around while being false of an operation introduced later. This is precisely the situation Section 7 identifies with the onset of repair.

\subsection{Clio as a Fixed Projection: The Admissibility Quotient}

The $L_2$-normalization instance above is the simplest possible member of a more general construction developed elsewhere in this research program, and stating that construction properly both sharpens the sufficiency claim and resolves an apparent inconsistency in how Clio has been formalized across different documents. Let $\mathrm{Adm}: X \to \mathcal{P}(X)$ assign to each state its admissible continuation set, and define the admissibility equivalence relation
\[
x \sim_{\mathcal{A}} y \iff \mathrm{Adm}(x) = \mathrm{Adm}(y).
\]
The \emph{admissibility quotient} $M_0 = X / {\sim_{\mathcal{A}}}$ comes with a canonical projection $\pi_{\mathcal{A}}: X \to M_0$, $\pi_{\mathcal{A}}(x) = [x]$, satisfying $\pi_{\mathcal{A}}(x) = \pi_{\mathcal{A}}(y) \iff \mathrm{Adm}(x) = \mathrm{Adm}(y)$ by construction. A projection $\pi: X \to M$ is \emph{exact minimal admissibility-preserving} if it is exact (no admissibility-relevant distinction is collapsed) and minimal (nothing beyond admissibility-class membership is retained). Under this definition, $\pi_{\mathcal{A}}$ satisfies a universal property at the algebraic level with no topology assumed: since $\pi$ is exact, it is constant on each admissibility class, and therefore factors as $\pi = \bar{\pi} \circ \pi_{\mathcal{A}}$ for a unique map $\bar{\pi}: M_0 \to M$, injective by minimality. Continuity of $\bar{\pi}$ is a further claim, not a free consequence of this factorization, and requires $X$ and $M$ to carry topologies under which $\sim_{\mathcal{A}}$ is a closed equivalence relation and $M_0$ is given the quotient topology, in which case $\bar\pi$'s continuity follows from the quotient topology's universal property in the usual way. This essay states the factorization at the algebraic level, which is all that is needed for what follows, and does not assume the additional topological structure needed for the continuous version.

The information-theoretic consequence needs a comparable qualification. Where $I(X;M)$ denotes mutual information under a specified probabilistic structure, equality $I(X;M) = I(X;M_0)$ for two exact minimal admissibility-preserving projections holds cleanly when the two projections differ only by a measurable bijective relabeling of $M_0$, under whatever measure-theoretic conditions make both mutual informations well-defined in the first place; for deterministic maps on continuous spaces without further assumptions, $I(X;M)$ can fail to be finite or well-defined at all, and this essay does not assume away that difficulty. What survives without extra hypotheses is the algebraic factorization above: any two exact minimal admissibility-preserving projections are related by an injective relabeling, and any further claim about equal information content is a claim about that relabeling's measure-theoretic behavior, to be checked case by case rather than assumed.

Two things are worth being precise about here, because an earlier draft of this result in the same research program stated it more strongly, as an $\argmin$ over all admissibility-preserving projections, and that stronger statement was subsequently identified as unsupported and withdrawn in favor of an equality-of-information-content statement. The corrected version proves that exact minimal admissibility-preserving projections agree in their information content, subject to the measure-theoretic qualification just given, not that $M_0$ minimizes complexity over some larger comparison class that includes non-exact or non-minimal projections; the two claims look similar but are not the same, and this essay uses the corrected, weaker one.

Under this general schema, the $L_2$-normalization instance $q_{\mathrm{Clio}}$ from the opening of this section is a specific candidate projection, not the definition of Clio itself. It is exact minimal admissibility-preserving exactly when the task's admissibility structure happens to be scale-invariant, which is precisely the condition the decisive inequality above already states in the vocabulary of Section 2; the general schema simply supplies the geometric object, $\mathrm{Adm}$ and its quotient, that the earlier condition was implicitly a condition about. For the record: this normalization instance is attested directly in an earlier essay outline in this program, on activation-space geometry, under the heading of $L_2$ normalization as Clio projection, magnitude discarded and direction preserved; it was not invented for this essay, but it is, exactly as that outline already treated it, one illustrative instance of the general projection $\pi: X \to M$ rather than a claim about Clio's canonical form. Nothing in the wider corpus establishes normalization as Clio's defining or most representative case, and this essay does not claim otherwise.

\subsection{Clio as Iterative Estimation of the Quotient}

A second formalization found elsewhere in this program does not describe $\pi_{\mathcal{A}}$ directly but instead an algorithm for estimating it from data. Define a regularized reduction functional over a hypothesis class $\Pi$ of candidate projections,
\[
L_{\mathrm{Clio}}[\pi] = L_{\mathrm{task}}(\pi) + \lambda\, \mathbb{E}_{x \sim \mathcal{C}}\bigl[E_\pi(x)\bigr] + \mu\, C(M),
\]
where $L_{\mathrm{task}}$ is task performance loss, $E_\pi(x)$ is projection error against the empirical admissibility structure estimated from a corpus $\mathcal{C}$, $C(M) = I(X;M)$ is representational complexity, and $\lambda, \mu > 0$ are regularization weights. Compactness of $\Pi$, uniform convergence of the empirical loss to its population counterpart, and a corpus $\mathcal{C}_n$ drawn i.i.d.\ with full support on $X$ are not by themselves sufficient to conclude that the minimizer converges to $\pi_{\mathcal{A}}$; they control the estimation problem but say nothing about which projection the population objective actually singles out. Two further ingredients are needed. First, the population version of $L_{\mathrm{Clio}}$ must have $\pi_{\mathcal{A}}$ as its unique minimizer, or as its unique minimizer modulo admissibility-preserving relabelings of the kind discussed above, since otherwise uniform convergence of the loss gives consistency toward whatever the population objective actually minimizes at, not toward $\pi_{\mathcal{A}}$ specifically. Second, $\Pi$ must be equipped with a metric under which both the stated compactness and uniform convergence hold and convergence of projections is itself measured, since "converges to $\pi_{\mathcal{A}}$" is not meaningful prior to fixing what topology that convergence is stated in. Under compactness of $\Pi$ in such a metric, uniform convergence of the empirical loss, a corpus $\mathcal{C}_n$ drawn i.i.d.\ with full support on $X$, and the identification condition that $\pi_{\mathcal{A}}$ (or its relabeling-equivalence class) is the unique population minimizer, the minimizer $\pi_n^{*}$ of $L_{\mathrm{Clio}}$ over $\mathcal{C}_n$ converges in probability to $\pi_{\mathcal{A}}$, in the fixed metric on $\Pi$, as $n \to \infty$.

This resolves what would otherwise look like a contradiction: is Clio a fixed geometric object or a learning algorithm. It is both, at two different levels, and the levels are not in tension once separated. $\pi_{\mathcal{A}}$ is the fixed target the construction of Section 4.1 defines exactly. $\pi_n^{*}$ is a data-driven estimator of that same target, consistent under the stated hypotheses but never guaranteed to equal $\pi_{\mathcal{A}}$ at any finite $n$. An $F$-sufficiency claim made about $\pi_n^{*}$ for finite $n$ is therefore always a task-relative claim in the sense of Section 2, carrying an estimation-error term in addition to whatever bias the projection's functional form contributes, even in the idealized case where $\pi_{\mathcal{A}}$ itself would be exactly sufficient.

A third, looser formalization appears in an earlier draft in the same program, describing Clio as a higher-order operator under which both the state and the projection mechanism co-evolve, $(X_t, \pi_t) \mapsto (X_{t+1}, \pi_{t+1})$, with no compactness or uniform-convergence hypotheses attached. This essay treats that formulation as superseded by the convergence result above, in the same spirit as the correction noted in Section 4.1: a co-evolving, unconstrained update rule is a strictly harder object to make sufficiency claims about than a convergent estimator of a fixed target, and nothing in the surviving formalization requires the mechanism itself to be non-stationary. Where a future version of Clio genuinely does require the projection mechanism to evolve rather than merely converge, that version would need its own convergence theory, stated with its own hypotheses, rather than inheriting the guarantees proved here.

\subsection{Collapse as a Quantitative Failure of Sufficiency}

The binary counterexample above, $f_{\mathrm{mag}}$ crossing a magnitude threshold, is an instance of a more general and independently formalized notion in this research program: a projection $\pi$ \emph{collapses} a distinction $(x_1, x_2)$ if $x_1 \neq x_2$ but $\pi(x_1) = \pi(x_2)$, and the collapsed distinction is \emph{reachability-relevant}, meaning some task-relevant reachability set differs between $x_1$ and $x_2$ even though $\pi$ cannot see the difference. This is exactly $F$-sufficiency failure restated in geometric language, with $F$ taken to be the class of functions sensitive to reachability from $x_1$ versus $x_2$; the magnitude-threshold function $f_{\mathrm{mag}}$ constructed above is a specific witness that this class is nonempty for $q_{\mathrm{Clio}}$.

Where the binary counterexample only shows that collapse can occur, a separately developed quantitative measure gives collapse a magnitude. Over the fiber $\pi^{-1}(m)$ above a point $m$ in the projected space, define a collapse observable
\[
\Lambda(m) = \mathbb{E}_{x_1, x_2 \sim \pi^{-1}(m)}\bigl[D_{\mathrm{KL}}(p_{x_1} \,\|\, p_{x_2})\bigr],
\]
the expected divergence between the task-relevant distributions associated with two points the projection has identified. $\Lambda(m) = 0$ exactly when the fiber over $m$ is task-homogeneous, and $\Lambda(m) > 0$ measures how badly the projection has conflated distinguishable, task-relevant states at $m$. The counterexample of the preceding subsection is the case $\Lambda(m) > 0$ for the fiber containing both $x$ and $y = \alpha x$ on opposite sides of $\tau$.

This gives the paper two candidate quantitative accounts of the same qualitative failure, an information-theoretic one via $C(M) = I(X;M)$ from Section 4.1 and a divergence-theoretic one via $\Lambda(m)$ here, and this essay does not attempt to reconcile them into a single measure. Both are internally coherent and both correctly diagnose the binary counterexample as a case of insufficiency; whether $\Lambda$ integrated over the fiber space stands in a precise quantitative relationship to $H(X) - H(M)$, the mutual-information gap the exact minimal projections of Section 4.1 are built around, is a genuine open technical question this essay leaves unresolved rather than answering by assertion.

\section{From Joint Addressing to Operational Coupling}

The JPEG and Clio cases both concern reconstruction accuracy: does $g(q(x))$ agree with $f(x)$, exactly or approximately. This section asks a different kind of question, one that does not reduce to reconstruction accuracy at all: given several signals that a system consumes jointly, do they merely happen to be retrieved together, or does the system's actual transition dynamics depend on their joint configuration in a way that cannot be pulled apart into independent contributions.

\subsection{The R/S byte as a controlled baseline}

The same JPEG decoding architecture supplies a useful baseline case for this question, distinct from the length-classification mechanism of Section 3. Once a Huffman code has been decoded, its associated symbol is a single byte encoding two logically distinct quantities: a run length $R$, the number of preceding zero-valued coefficients, and a size category $S$, the number of additional bits needed to specify a nonzero coefficient's value. These two quantities are packed into one byte, with $R$ occupying the upper nibble and $S$ the lower.

This is joint addressing, not operational coupling. $R$ and $S$ are stored and retrieved together because every decoding step needs both, but each remains independently recoverable by simple masking and shifting: $R = \mathrm{byte} \gg 4$ and $S = \mathrm{byte} \,\&\, \mathtt{0xF}$. Nothing about the decoder's downstream behavior given $R$ depends on the particular value of $S$, or vice versa; the packing is purely a storage and addressing optimization; the two fields could be unpacked into separate bytes with no change to the decoder's logic beyond an extra step. This gives the middle of the hierarchy in Section 6 a concrete, minimal instance: co-location plus independent recoverability, with no interaction in how the two fields are subsequently used.

\subsection{Hydra's transition operator}

Hydra's cognitive architecture combines a perception signal $P_t$ (from PERSCEN), a relevance signal $R_t$ (from Relevance Activation Theory), and a memory signal $M_t$ (from Chain of Memory) into a state transition
\[
z_{t+1} = \mathcal{H}(P_t, R_t, M_t; \Phi_t, \mathbf{v}_t, S_t),
\]
where $\Phi_t$, $\mathbf{v}_t$, and $S_t$ denote the ambient RSVP scalar, vector, and entropy fields at time $t$. The question this section poses is whether $P$, $R$, and $M$ enter $\mathcal{H}$ the way $R$ and $S$ enter the JPEG symbol byte, as jointly retrieved but independently substitutable quantities, or whether they enter it the way a genuinely coupled system's variables enter a nonlinear transition, unable to be varied one at a time without changing what the others mean.

A first, tempting formalization of separability is vacuous and must be rejected before any real test can be built. Stating that $\mathcal{H}$ is separable whenever
\[
\mathcal{H}(P, R, M) = h\bigl(h_P(P), h_R(R), h_M(M)\bigr)
\]
for \emph{some} functions $h, h_P, h_R, h_M$ is trivially satisfiable by every transition operator whatsoever: take $h_P, h_R, h_M$ to be the identity and $h = \mathcal{H}$ itself. This formalization asserts nothing, because it places no restriction on the composition class $h$ is drawn from.

A non-vacuous version requires declaring the composition class in advance. Fix a restricted family $\mathcal{F}_{\mathrm{sep}}$ of candidate separable approximations, for instance the additively separable family
\[
\widetilde{\mathcal{H}}(P, R, M) = h_P(P) + h_R(R) + h_M(M)
\]
for some fixed link function, or a comparably restricted multiplicative or low-rank family appropriate to the signals involved. Separability under $\mathcal{F}_{\mathrm{sep}}$ then holds when
\[
\inf_{\widetilde{\mathcal{H}} \in \mathcal{F}_{\mathrm{sep}}} \mathbb{E}\left[ L\bigl( \mathcal{H}(P,R,M), \widetilde{\mathcal{H}}(P,R,M) \bigr) \right] \leq \varepsilon_{\mathrm{sep}}
\]
for a stated loss $L$ and threshold $\varepsilon_{\mathrm{sep}}$. Operational coupling, on this formalization, is precisely the residual interaction that no member of the declared family $\mathcal{F}_{\mathrm{sep}}$ can absorb. This makes the coupling claim relative to a stated composition class rather than absolute, which is the correct level of honesty for the claim: a system might resist an additively separable approximation while admitting a different, still-separable, multiplicative one, and the paper's claim about Hydra must specify which family was tested and failed.

Where $\mathcal{H}$ is differentiable, a local diagnostic supplements the global test above: nonzero mixed partial derivatives,
\[
\frac{\partial^2 \mathcal{H}}{\partial P\, \partial R} \neq 0, \qquad
\frac{\partial^2 \mathcal{H}}{\partial R\, \partial M} \neq 0, \qquad
\frac{\partial^2 \mathcal{H}}{\partial P\, \partial M} \neq 0,
\]
indicate that the marginal effect of one variable depends on the current value of another, which is direct evidence of interaction under the chosen coordinates. This diagnostic is coordinate-sensitive, since a reparameterization can in general remove or introduce nonzero mixed partials without changing the underlying system, and it establishes only that interaction exists under the coordinates tested, not that the interaction is functionally meaningful or exploited by the system's actual behavior.

The strongest test available is empirical rather than analytic: intervention. Fix two of the three signals and vary the third, then ask whether the resulting transition $z_{t+1}$ remains admissible and interpretable under the system's own admissibility criteria. If independent interventions on $P$, $R$, or $M$ routinely destroy the meaning of the remaining variables, that is, if holding $R$ and $M$ fixed while varying $P$ produces transitions the system itself cannot classify as coherent continuations, then Hydra is operating beyond joint addressing in a sense the ablation and mixed-partial tests can only gesture toward. This essay does not claim to have run this intervention test on Hydra; it specifies the test as the criterion a future implementation-level paper would need to satisfy before asserting operational coupling as more than a suggestive architectural resemblance.

A relevant architectural predecessor for treating perception, judgment, and decision-making as an integrated cognitive system rather than a loose federation of independently invoked modules is de Gyurky and Tarbell's design for autonomous deep-space spacecraft cognition, which draws explicitly on Kant's account of judgment and reason and frames integrated autonomous cognition as what the authors call a modern categorical imperative \cite{degyurky2014}. This source does not establish that Hydra satisfies operational coupling in the formal sense defined above; no existing published account of Hydra has been tested against the separability family or intervention criteria given here. What it supplies is a historically grounded example of a system explicitly designed around the premise that perception, judgment, and decision cannot be cleanly separated into independently retrievable modules, against which Hydra's stronger, formally testable separability claim can be stated and eventually checked.

\subsection{A Complication for the Intervention Test: Memory as Stabilized Residue}

The intervention criterion proposed above, holding two of $P$, $R$, $M$ fixed while varying the third, presupposes that a signal can be held fixed by simply not querying it during the test. A result developed elsewhere in this program for Hydra's memory component specifically shows this presupposition can fail. In a conventional vector database, retrieval and storage are distinct operations: $\mathrm{retrieve} \neq \mathrm{store}$, and querying an entry leaves it unchanged. In a residue-field memory architecture of the kind Hydra's Chain of Memory component is built on, this separation does not hold: $\mathrm{retrieve} \approx \mathrm{reinforce}$, because every traversal of the memory field deepens the trace it follows, and reading a memory state is itself a write to that state.

If $M_t$ is a residue field in this sense, then \emph{holding $M$ fixed} while an intervention test varies $P$ or $R$ is not experimentally well-defined in the naive sense the test assumed: consulting $M_t$ at all, even passively, as part of evaluating $\mathcal{H}(P_t, R_t, M_t; \ldots)$ under the intervention, changes $M_t$ into $M_t'$ before the next step of the test can be run, so the "fixed" component has already moved. This is not a flaw in the separability question itself, but it means the intervention criterion of Section 5.2 needs a more careful protocol than simply not touching two of the three signals. A workable version would perform an external substitution, replacing $M_t$ with a counterfactual $M_t'$ constructed outside the system's own read/traversal dynamics, prior to the intervention, rather than attempting to query $M_t$ in place and treat the query as neutral. Whether such an external substitution is itself faithful to what "holding memory fixed" should mean for a residue-field architecture is a design question this essay raises but does not resolve; it is flagged here because it is a genuine complication that a system built on a passive-readout memory model, such as the R/S byte's storage, simply does not have.

\subsection{Mixture-of-Experts as a Worked Separability Contrast}

It is worth having one worked example of an architecture that looks coupled on inspection but tests separable under a properly declared family, so that Hydra's open status is not mistaken for a default answer. A mixture-of-experts architecture combines several component outputs $h_1(x), \ldots, h_k(x)$ through a learned gating function into a single result,
\[
\mathcal{H}_{\mathrm{MoE}}(x) = \sum_{i=1}^{k} w_i(x)\, h_i(x), \qquad w_i(x) = \mathrm{gate}_i(x),
\]
and superficially resembles operational coupling: the final output depends jointly on every expert and on the gate, and no single expert's contribution is visible in isolation from the final sum. Under the declared family $\mathcal{F}_{\mathrm{sep}}$ consisting of gated weighted sums of independently computed component functions, exactly the form above, $\mathcal{H}_{\mathrm{MoE}}$ is a member of $\mathcal{F}_{\mathrm{sep}}$ by construction, so the residual-interaction test of Section 5.2 returns zero: there is no interaction left over that a member of the declared family fails to absorb, because the architecture's own combination rule already is such a member. This establishes separability under the declared family, not joint addressing specifically; the experts $h_1, \ldots, h_k$ can in principle be computed and stored entirely independently before their outputs are aggregated, so mixture-of-experts is better described as a separable compositional architecture than as an instance of joint addressing in the sense the R/S byte instantiates it. What it does share with the R/S byte is the more general point the hierarchy above is built around: nonlinear, jointly consumed combination is not by itself evidence of operational coupling, since a properly declared family can absorb it cleanly. This is a genuine contrast with Hydra's open case: Hydra's status is unresolved not because joint, nonlinear combination is automatically evidence of coupling, since the mixture-of-experts case shows it need not be, but because no comparably explicit gated-sum family has yet been shown to absorb $\mathcal{H}$'s interaction among $P$, $R$, and $M$, and the preceding subsection's memory-reinforcement complication means even constructing the test data to check this is not as simple as running the mixture-of-experts case's ablations directly.

\subsection{The Named Layer Equations, and Where Coupling Might Actually Live}

The abstract transition $z_{t+1} = \mathcal{H}(P_t, R_t, M_t; \Phi_t, \mathbf{v}_t, S_t)$ used above is a deliberate simplification of a more specific architecture recorded elsewhere in this program, and making that architecture explicit sharpens rather than settles the separability question. Hydra's relevance layer, RAT, is specified as a cue-density field
\[
\rho_c(x) = \exp\left[-\tfrac{1}{2}(x - \mu_c)^{T} \Sigma_c^{-1} (x - \mu_c)\right],
\]
with relevance-directed movement given by gradient ascent on that field, $\dot{x} = \nabla \rho_c(x)$. This is significant for the separability question because a gradient-flow contribution is structurally additive: if $R_t$'s entire effect on the transition is to add a term $\nabla \rho_c(x)$ to an otherwise independent dynamics, that particular contribution belongs to the declared additive family $\mathcal{F}_{\mathrm{sep}}$ of Section 5.2 by construction, and no separate residual-interaction test is needed to place it there.

The coupling this observation cannot rule out is coupling through the parameters $\mu_c, \Sigma_c$ themselves. If the cue center $\mu_c$ or covariance $\Sigma_c$ is derived from the current perception state $P_t$ or memory state $M_t$, rather than fixed externally, then $R_t$'s contribution is additive in its functional form while still coupled to $P$ and $M$ through its parameterization, and the residual-interaction test of Section 5.2 must be applied to the parameter-dependence $\mu_c(P_t, M_t), \Sigma_c(P_t, M_t)$ rather than to the top-level sum, or it will register a false negative: a system can look separable at the level of the combination rule while smuggling the actual coupling into how each term's parameters are computed. This is worth stating as a general caution rather than a fact specific to RAT:

\begin{proposition}[Parameter-Channel Coupling]
Let a transition take the displayed additive form
\[
\dot{x} = F(x; P, M) + \nabla \rho_c\bigl(x;\, \mu_c(P, M),\, \Sigma_c(P, M)\bigr)
\]
for some fixed functional forms $F$ and $\rho_c$. Membership of this expression in an additively separable family $\mathcal{F}_{\mathrm{sep}}$, tested at the level of the displayed sum, does not imply that $R$'s contribution is independent of $P$ and $M$, whenever the parameters $\mu_c, \Sigma_c$ are themselves nonconstant functions of $P$ and $M$. Top-level additive form is a claim about the combination rule; independence of the components is a separate claim about the arguments each term is allowed to depend on, and the residual-interaction test of Section 5.2 only detects the latter if it is applied to the full parameter dependence rather than to the displayed functional skeleton.
\end{proposition}

This is probably the single most useful mechanism-level observation available for auditing an architecture like Hydra in practice: a system audited only at the level of its displayed equations, checking that a sum of terms looks additively separable, can pass that check while still routing genuine coupling through the arguments each term is fed, and the audit has to follow the parameters, not just the combination rule, to be trustworthy.

The same caution applies to Hydra's memory layer, CoM, whose recurrence
\[
M_{i+1} = \phi(M_i, u_i, c_i)
\]
depends on a contextual input $c_i$ that Section 5.1's residue-field result already complicates; if $c_i$ is itself derived from $R_t$'s relevance field or $P_t$'s perception state, the same parameter-channel coupling applies to memory's update as to relevance's gradient.

The full pipeline is recorded as a sequential composition of operators rather than a single joint function of three simultaneous arguments,
\[
\mathcal{H} = \mathrm{GLU}_{\mathrm{RSVP}} \circ M \circ T \circ F_a \circ G_a \circ \rho_c,
\]
with cue activation feeding personalized graph transformation, tiling, memory evolution, and RSVP-gated gluing in that order. This composition structure is itself informative and was obscured by treating $P$, $R$, $M$ as co-arguments to one operator: a sequential pipeline is not the same separability question as a joint function of simultaneous inputs, and the more precise question is whether each stage's output can be perturbed independently of the stage before it without changing what the later stages do with it, evaluated one adjacent pair of stages at a time, rather than asking whether $P$, $R$, and $M$ jointly determine $\mathcal{H}$ all at once. This essay does not carry out that stage-by-stage analysis; it identifies the composition structure as the correct object for a future analysis to test, in place of the flattened three-argument transition used earlier in this essay for expository purposes.

A further diagnostic, structurally independent of the residual-interaction test, the mixed-partial test, and the intervention test of Section 5.2, is available if Hydra's assembly of local structures is treated categorically as a colimit or gluing construction, $\operatorname{colim}_{\mathcal{C}} \mathcal{A}$, over local cue states, graph sections, tiles, and memory trajectories. Under this reading, a failure to assemble the components consistently appears not as a low score but as a gluing obstruction $\omega$, formally a class in the relevant cohomology witnessing that the local sections cannot be glued under the proposed compatibility maps; $\omega = 0$ exactly when a global gluing exists. This diagnostic must be read carefully, since it is easy to overstate what it shows. $\omega \neq 0$ proves that the attempted gluing fails, that the proposed compatibility conditions on the local pieces are mutually incompatible; it does not by itself prove operational coupling in the transition-level sense defined above, and it says nothing about a transition that is in fact successfully computed and running, since a nonzero obstruction is a statement about a failed assembly, not a property of a working one. What it can supply, when $\omega \neq 0$, is a location: it identifies which local sections are incompatible and where the proposed gluing conditions break down, which is different information from, and does not substitute for, evidence that a successfully operating $\mathcal{H}$ resists the separable family $\mathcal{F}_{\mathrm{sep}}$. Whether Hydra's actual pipeline has $\omega \neq 0$ under any specific compatibility maps is, like the stage-by-stage separability question above, not established by the architectural description alone and remains open.

The relaxation dynamics governing the RSVP fields, $\Phi$, $\mathbf{v}$, $S$, are recorded elsewhere in the form $\dot{S} = -\gamma \int_{\Omega} \lVert \nabla S \rVert^2 \, dx$ under stated no-forcing, no-flux conditions. Taken literally this equates the time derivative of a field $S(x,t)$ with a single global scalar, which is not well-typed; the equation is better read as describing the derivative of an entropy functional $\mathscr{S}[S](t) = \int_\Omega S(x,t)\,dx$ or a comparable scalar summary of the field, giving
\[
\frac{d\mathscr{S}}{dt} = -\gamma \int_\Omega \lVert \nabla S(x,t) \rVert^2 \, dx,
\]
with the local field $S(x,t)$ kept notationally distinct from the global functional $\mathscr{S}[S]$ it feeds into. Where the source material states the relation in the original $\dot{S}$ form, this essay preserves that as a quoted structural relation rather than silently correcting it, but treats the functional form above as the one actually being asserted. Either way, the resemblance to diffusive relaxation toward a uniform field is architectural, not physical: whether these correspondences reduce to conventional thermodynamic, fluid, or gravitational equations is recorded elsewhere in this program as an explicitly open problem, and this essay relies on the relaxation dynamics only as a structural description of how the RSVP-gated final stage of $\mathcal{H}$ behaves, not as a claim that Hydra instantiates literal physical thermodynamics.



\section{A Hierarchy of Representational Integration}

The cases above support a three-level hierarchy of increasing integration, though the relationship between the levels is a claim about relative strength of assertion, not a set-inclusion relationship among literal implementation properties. Jointly addressed values can remain physically distributed across separate storage, and operationally coupled variables need not share a storage region or a retrieval address at all; the three properties are not automatically nested, and stating them with $\subsetneq$ without qualification would overclaim a structural relationship the definitions do not actually establish.

\emph{Shared representation}, or co-location, means only that several values occupy a shared storage region, with no claim about how they are retrieved or used. \emph{Joint retrieval}, or joint addressing, means one retrieval operation returns the values together, as in the R/S byte, while each value remains independently recoverable and independently consumable by whatever uses it afterward. \emph{Transition-level interdependence}, or operational coupling, means the system's transition semantics depend on the joint configuration of the values in a way that resists approximation within a declared separable composition class $\mathcal{F}_{\mathrm{sep}}$, as tested via the residual-interaction criterion, the mixed-partial diagnostic, or the stronger intervention criterion of Section 5. The two names given for each level are used interchangeably in what follows; the first emphasizes the implementation property, the second the strength of integration being claimed.

What can be stated correctly is that each level asserts a strictly stronger claim about integration than the one before it:
\[
\text{co-location} \;\prec\; \text{joint addressing} \;\prec\; \text{operational coupling},
\]
where $\prec$ denotes ``asserts a strictly stronger form of integration than,'' not set inclusion. Within the restricted class of architectures this essay has actually examined, ones where the stronger levels are built by adding integration on top of an already-present weaker level rather than replacing it, the inclusion reading and the stronger-assertion reading coincide: the R/S byte's joint addressing is built on top of ordinary co-located storage, and a hypothetical operationally coupled version of it would still be retrieved jointly. This essay does not claim the coincidence holds outside that restricted class, and a general architecture satisfying a stronger level while failing a nominally weaker one, coupled variables stored and retrieved entirely separately, for instance, is not excluded by anything argued here.

The JPEG R/S byte instantiates the first step cleanly: it is joint addressing and demonstrably not operational coupling, since $R$ and $S$ are recoverable by masking and shifting with no cross-dependence in how either is subsequently used. Whether Hydra instantiates the second step is the open, formally specified question Section 5 leaves for future empirical work.

\section{Discussion: Sufficiency, Admissibility, and Repair}

Compression creates an epistemic obligation rather than discharging one. A quotient $q$ that is $F$-sufficient for some function class $\mathcal{F}$ has not established that it is sufficient in general; it has established a scope, the class $\mathcal{F}$, within which its discarded distinctions are guaranteed not to matter. Section 3.1 already made this precise for the JPEG case: total and finitely checkable relative to one declared function, not relative to every function the same bits could support. Clio's scope is the same kind of bounded guarantee stated for a harder case: a standing assumption about which function class future downstream operations will belong to, an assumption that holds until some operation genuinely needs the discarded magnitude, at which point the projection's sufficiency claim, true of every task it was actually designed around, turns out to be false of the new one. Hydra's case adds a further dimension: the question is not only whether a discarded distinction will later be needed, but whether the components that appear jointly addressed were ever separable at all, so that the very notion of "which distinction was discarded" may not apply cleanly to a coupled system in the way it applies to a projection.

This is the precise sense in which sufficiency and repair are the same phenomenon viewed from two moments in time. A representation's sufficiency is a claim about which futures it is prepared to support; repair is what becomes necessary exactly when a future arrives that the representation's designer did not include in that claim. Outside the finite, fully enumerable case that JPEG happens to occupy, sufficiency is never a static, once-and-for-all property of a representation. It is a claim, always relative to a declared function class, always capable of being falsified by a task the class did not anticipate, and always, when it fails, pointing to exactly the distinction that must be recovered, reconstructed, or reintroduced before the system can proceed.

\subsection{Failure as Diagnostic, Not Noise}

A further result recorded elsewhere in this program complicates the convergence theorem of Section 4.2 in a way worth making explicit rather than leaving as an implicit tension. That theorem states that an iterative Clio estimator $\pi_n^{*}$ converges to the fixed admissibility quotient $\pi_{\mathcal{A}}$ under stated hypotheses; convergence is treated there as the success condition. A separately recorded analysis of projection failure argues that this framing can be exactly backwards in a specific circumstance: a projection that fails to converge, oscillating rather than settling, may be preserving a reachability-relevant distinction that a successfully converged, smoothed quotient would have erased. On this reading, non-convergence is not automatically evidence of a poorly behaved estimator; it can instead be evidence that the target $\pi_{\mathcal{A}}$ the estimator is being pushed toward is itself an inadmissible quotient, one that identifies states whose futures genuinely differ.

A recorded empirical instance of exactly this pattern involved persistent obstruction producing bistable or low-period oscillation under repeated projection and repair, in a coupled TARTAN-tiling and Clio-projection setting. The oscillation was diagnosed as structural non-gluing rather than ordinary numerical non-convergence, and the standard repair loop, iterating the projection to drive down apparent error, was found to over-smooth the state: each repair step reduced the visible discrepancy while erasing the very distinction whose persistence would have revealed that the quotient was too coarse to begin with. This suggests a genuine hazard, though it is a hazard of interpretation rather than a flaw in Section 4.2's theorem itself: that theorem, under its stated hypotheses, is a consistency result, and if $\pi_{\mathcal{A}}$ genuinely is the population minimizer identified there, the estimator $\pi_n^{*}$ can converge to it perfectly regardless of whether $\pi_{\mathcal{A}}$ happens to be a well-specified target for the task at hand. Statistical consistency establishes that the estimator is finding what it was built to find; it does not establish that what it was built to find is admissible. A more careful protocol would distinguish a genuine \emph{repair} regime, where continued iteration toward $\pi_{\mathcal{A}}$ is warranted because the discrepancy really is noise, from an \emph{expansion} regime, where the correct response to persistent, structured oscillation is to enlarge the representational space $M$ rather than to keep projecting harder onto the one already chosen. Section 4.2's theorem remains a correct consistency result for estimation of the stipulated target under either regime; what is conditional on the regime is not the theorem but the further, separate judgment that convergence constitutes epistemic success, which is warranted only once the target quotient has been established to belong to a repair regime rather than to require representational expansion.

\subsection{Clio and Hydra as a Composed Pipeline}

The two case studies of this essay are frequently treated as unrelated, but a compact formulation elsewhere in this program states their relationship precisely: Clio asks what a projection loses, and Hydra asks what an assembly creates. The two questions compose, rather than merely coexist, into a single pipeline,
\[
X \xrightarrow{\;\mathrm{Clio}\;} M \xrightarrow{\;\mathrm{Hydra\ gluing}\;} Z \xrightarrow{\;\mathrm{RSVP\ transition}\;} Z',
\]
in which Clio's quotient must justify the distinctions it discards before Hydra's assembly ever sees $M$, and Hydra's gluing must justify the interactions it introduces among the assembled components before the RSVP-gated transition acts on the result. Repair, in this composed picture, is not a single failure mode but has two structurally distinct entry points: it is needed at the Clio stage whenever the quotient turns out not to be $F$-sufficient for a function class a later task actually belongs to, in the sense of Sections 2 and 4, and it is needed at the Hydra stage whenever locally coherent components fail to glue into a consistent global state, in the sense of the obstruction diagnostic introduced above. Distinguishing which stage a given failure originates in matters practically: a $\pi_{\mathcal{A}}$ that is already exact minimal admissibility-preserving is not the right target for further repair effort if the actual defect is a nonzero gluing obstruction downstream at the Hydra stage, and conversely no amount of re-assembling already-admissible components will fix a quotient that discarded a distinction some function in the required class still needs.

\section{Conclusion}

Three architectures, examined through a single formal lens, occupy three distinct positions on a spectrum of sufficiency guarantees. JPEG's Huffman length classification achieves exact, finitely verifiable sufficiency because its domain is small enough to check exhaustively. Clio's normalization achieves, at best, sufficiency relative to a declared function class, a claim that must be argued and can be falsified by a single counterexample of the kind constructed in Section 4. Hydra's transition operator raises a question neither of the other two cases poses: whether its jointly consumed signals are separable at all, a question this essay has shown cannot be answered by an unrestricted factorization claim, and has instead reformulated as a test against a declared composition class, supplemented by a coordinate-sensitive local diagnostic and a stronger empirical intervention criterion. What unifies the three cases is not that they compress by the same mechanism, but that each makes a claim, of differing strength and differing verifiability, about which distinctions a downstream operation will need, and each therefore incurs an obligation whose failure mode is exactly what the admissibility and repair vocabulary used elsewhere in this research program was built to describe.

\begin{thebibliography}{9}

\bibitem{mitchell2002}
Mitchell, J.~L., Cazes, A.~N., and Leeder, N.~M. (2002). Fast JPEG Huffman Encoding and Decoding. U.S. Patent No. 6,373,412 B1.

\bibitem{degyurky2014}
de Gyurky, S.~M. and Tarbell, M.~A. (2014). \textit{The Autonomous System: A Foundational Synthesis of the Sciences of the Mind}. Hoboken, NJ: John Wiley \& Sons.

\end{thebibliography}

\end{document}
